\documentclass[10pt,a4paper]{amsart}
\usepackage[margin=19mm]{geometry}
\usepackage{amsmath,amssymb,mathtools}
\usepackage[scr=rsfs,cal=euler]{mathalfa}
\usepackage{xfrac}
\usepackage[unicode,hidelinks,bookmarks=false]{hyperref}
\newcommand{\ie}{i.e.\ }
\newcommand{\cf}{cf.\ }
\newcommand{\resp}{respectively\ }
\newcommand{\Subsection}{\subsection}
\providecommand{\nobreakdash}{\nobreak\hbox{-}}
\setlength{\emergencystretch}{2em}
\multlinegap0pt
\usepackage{amssymb}
\usepackage{mathtools}
\usepackage[shortlabels]{enumitem}
\newtheorem{proposition}{Proposition}
\newcommand{\Alt}{\operatorname{Alt}}
\newcommand{\FSym}{\operatorname{FSym}}
\newcommand{\N}{\mathbb N}
\newcommand{\Ra}{\mathcal{R}}
\newcommand{\Z}{\mathbb Z}
\renewcommand{\ker}{\operatorname{Ker}}
\newcommand{\supp}{\operatorname{supp}}
\renewcommand{\wr}{\operatorname{\,wr\,}}
\title{Finiteness properties of Subgroups of Houghton Groups of full Hirsch length}
\author{Charles Garnet Cox, Peter Kropholler, Armando Martino}
\date{Source: arXiv:2508.07816v3 (CC BY 4.0)}
\begin{document}
\maketitle

\noindent\emph{Source and licence.} Finiteness properties of Subgroups of Houghton Groups of full Hirsch length. Version arXiv:2508.07816v3 (CC BY 4.0), distributed by arXiv under the Creative Commons Attribution 4.0 International licence, http://creativecommons.org/licenses/by/4.0/, as recorded in the arXiv licence metadata for this exact version and shown on the abstract page https://arxiv.org/abs/2508.07816v3. Full licence notice: \texttt{licenses/arxiv-2508.07816-CC-BY-4.0.txt}.\par\smallskip

\noindent\textit{Editorial scope.} Counted unit: the article's proposition H\_2 (source0.tex lines 1430--1440). The article's complete original proof of it is delivered with it (source0.tex lines 1441--1504, one \texttt{proof} environment of 64 lines). No mathematics is written, repaired or re-derived: the statement and the proof are the article's own lines, in the article's own notation, and no retained line is substituted or re-flowed.  No further source-local context is needed: every cross-reference of every retained line resolves inside the two delivered spans.  Nothing was omitted from any retained span, so the package carries no omission record.  No retained line contains a citation, so the wrapper carries no bibliography and no bibliography entry.  No retained line contains a figure environment, an image-inclusion command or drawing code, so no image is delivered and the package declares no asset dependency.  Editorial formatting only, in addition: an AMS \texttt{amsart} preamble that restates the article's own declarations and macros used by the retained lines, namely the environments \texttt{proposition} on separate counters, together with the standard packages those lines need.  The retained environments print in this document as Proposition 1 in order of appearance, whereas the article prints the same items with the numbering of its own full document.  The complete native source of the article as distributed in the arXiv e-print for this exact version is bundled unchanged as \texttt{sources/183053/source0.tex}.
\begin{proposition}
	\label{normal subgroups of wreaths in H_2}
	Consider the permutational restricted wreath product, $A \wr_{\Ra_n} \Gamma$, where $\Gamma$ is either a finite index subgroup of $H_n$ or $n=2$ and $\Gamma \cong \Z$ acting on $\Ra_2$ by translations and $A$ is some finite group.
	
	
	Let $\pi: A \wr_{\Ra_n} \Gamma \to \Gamma$ denote the natural projection map and let $G$ be a subgroup of $A \wr_{\Ra_n} \Gamma$ such that $\pi(G) = \Gamma$. Then 
	\begin{itemize}
		\item $G$ is finitely generated, and
		\item $G$ has \textbf{max-n}. 
	\end{itemize}
\end{proposition}

\begin{proof} 
	
We will start the proof by showing that any normal subgroup $N$ of $G$ which is contained in $G \cap B $ is finitely normally generated in $G$. 

	
	Let $B=\ker \pi$, the base of the wreath product. Then $B = \bigoplus_{\Ra_n} A$. Since this is a direct sum, every element of $B$ has a support which is a finite subset of $\Ra_n$. We also have the projection maps, $\rho_x: B \to A$ for every $x \in \Ra_n$.


	We deal with the case where $\Gamma$ is a finite index subgroup of $H_n$. 


To do this, first pick a basepoint $* \in \Ra_n$ and consider a finite subset, $F_0$ of $N$ such that $\rho_*(F_0) = \rho_*(N)$. (This is clearly possible, as $A$ is finite). 
The crucial property we will use is that $\Gamma$ contains $\Alt(\Ra_n)$ and hence acts $k$-transitively on $\Ra_n$ for any $k \in \N$. 


We let $S = \bigcup_{f \in F_0} \supp(f)$ be the finite subset of $\Ra_n$ consisting of the union of all the supports of elements of $F_0$. Then let $F$ denote the finite subset of $N$ whose supports lie in $S$. 

We note that since $\Gamma$ acts transitively on $\Ra_n$ and $\pi(G) = \Gamma$, we get that for any $x \in \Ra_n$ and any $n \in N$, there exists a $g \in G$ and an $ f \in F$ such that $\rho_x(f^g) = \rho_x(n)$. In fact, we could have taken the $f \in F_0$. However, do be aware that in general, $\rho_*(f) \neq \rho_x(f^g)$, even when $xg=*$, although they are conjugate in $A$ in that case. However, $|\rho_x(N)|$ is constant as we vary $x$, since $N$ is normal and $G$ acts transitively on $\Ra_n$. Therefore the elements $\rho_x(f^g)$ exhaust $\rho_x(N)$.  

We claim that $F$ normally generates $N$. We argue by contradiction. If this is not the case, choose some $n \in N$ outside of the normal closure of $F$ whose support is minimal. There are two cases to deal with. If the support of $n$ is not larger than $|S|$ then, since $\Gamma$ acts in a highly transitive way, there exists a $g \in G$ such that $n^g$ has support contained in $S$. Hence $n^g \in F$ or, equivalently, $n \in F^{g^{-1}}$, so that $n$ lies in the normal closure of $F$ resulting in a contradiction. Otherwise, the support of $n$ has at least $|S|$ elements. Choose some $x \in \Ra_n$ which lies in the support of $n$. 


Now, since the action of $\Gamma$ is highly transitive, we can find a $g \in G$ and a $f \in F$ satisfying the following: 
\begin{itemize}
	\item $\rho_x(f^g) = \rho_x(n)$
	\item The support of $f^g$ is contained in the support of $n$. 
\end{itemize}
But now, as $n$ does not lie in the normal closure of $F$, neither does ${(f^g)}^{-1} n$. But this element has a smaller support by construction. So the resulting contradiction shows that $N$ is finitely normally generated. Thus we have proved that any normal subgroup of $B$ is finitely normally generated in the case where $\Gamma$ has finite index in $H_n$. 

\medspace

We next prove the same property when $\Gamma=\Z$ and $n=2$. The proof is similar but we need to modify some details as follows. The set $F_0$ is defined in the same way. Now, since $\Ra_2 = \{ 1,2\} \times \N$, we will identify $\Ra_2$ with $\Z$ (this is not an order preserving identification, but it doesn't matter for what follows). Then the action of $\Z$ on $\Ra_2$ is simply the regular action of $\Z$ on itself. Given any finite subset $S$ of $\Z$, we then assign an interval, $I(S):=[\inf(S),\sup(S)]_{\Z}$ and the length of $S$, $l(S)$ to be $|\sup(S) - \inf(S)|$. 

Note that for any $b \in B$, its support has a length and this length is a conjugacy invariant. 

Now let $L$ be the maximum of the lengths of elements of $F_0$ and extend $F_0$ to a larger finite set, $F$, such that any $n \in N$ of length $L$ is conjugate to an element of $F$. As before, we claim that $F$ normally generates $N$ in $G$. This follows as before. We argue by contradiction, choosing an $n \in N$ of minimal length which does not lie in the normal closure of $F$. By construction, as before, the length of $n$ must be greater than $L$ (as otherwise it would be equal to a conjugate of an element of $F$). Then, as before, we can find an $x$ in the support of $n$, a $g \in G$ and a $f \in F$ such that: 
\begin{itemize}
	\item $\rho_x(f^g) = \rho_x(n)$
	\item The interval of the support of $f^g$ is contained in the interval of the support of $n$. 
\end{itemize}
It is then clear that the length of the support of ${(f^g)}^{-1} n$ is smaller than the length of $n$; this contradiction proves that $F$ normally generates $N$.

\medspace

Thus we have shown in either case that any normal subgroup of $G \cap B$ is finitely normally generated in $G$. We now use this to get the remaining results. 


\medspace

We next prove that $G$ is finitely generated. Since $G \cap B$ is a normal subgroup of $G$, it is normally generated by a finite set, $F$. Moreover by possibly enlarging $F$, we assert that there is a finite subset $S$ of $\Ra_n$ so that:
\begin{itemize}
	\item Any $f \in F$ has support contained in $S$ and, 
	\item Any element whose support is contained in $S$ is in $F$. 
\end{itemize}
That is, $F$ consists of exactly those (finitely many) elements whose support is contained in a fixed finite set, $S$. (We are using the fact that $A$ is finite here.)

Now since $\Gamma$ is finitely generated and $\pi(G)= \Gamma$ we can find a finite subset $X$ of $G$ such that $\pi(\langle X \rangle) = \Gamma$. That is, we take pre-images of a finite generating set. We claim that $G$ is generated by $X$ and $F$. Since $F$ normally generates $G \cap B$ and $\pi(\langle X \rangle) = \Gamma$, it is sufficient to show that $F^g \subseteq \langle X, F \rangle$ for all $g \in G$. Given such a $g$, there exists a $w \in \langle X \rangle$ such that $\pi(g)=\pi(w)$. But then the support of $F^{g w^{-1}}$ equals the support of $F$ (since $g w^{-1}$ acts trivially on $\Ra_n$). Hence $F^{g w^{-1}}  = F$, by the comments above (and the fact that $F$ is finite). In particular, $F^g = F^w \subseteq \langle X, F \rangle$. Hence $G$ is generated by $X$ and $F$. This shows that $G$ is finitely generated. 

\medspace

Finally, we prove that an arbitrary normal subgroup $N$ of $G$ is finitely normally generated. We have already proved that $N \cap B$ is finitely normally generated. And since every normal subgroup of $\Gamma$ is finitely normally generated, it follows that $NB/B \leq GB/B \cong \Gamma$ is also finitely normally generated, since $NB/B$ is a normal subgroup of $GB/B$. (Note when $\Gamma$ is a finite index subgroup of $H_n$, then a normal subgroup is either trivial, $\Alt, \FSym$ or finite index in $H_n$ and hence is finitely normally generated. When $\Gamma=\Z$, then every subgroup is finitely generated.)

It then follows that there is a finite subset $F_1$ of $N$ such that $F_1B$ normally generates $NB/B$ and a finite subset, $F_2$ of $N \cap B$ which normally generates $N \cap B$. Hence $N$ is finitely normally generated by $F_1 \cup F_2$. 
\end{proof}

\end{document}
